\documentclass[12pt,a4paper]{article}
\usepackage[UTF8]{ctex}
\usepackage{amsmath,amssymb,amsthm}
\usepackage{geometry}
\usepackage{graphicx}
\usepackage{hyperref}
\usepackage{bm}

\geometry{margin=2.5cm}

\title{$\dot{T} = [V]T$的详细证明}
\author{Modern Robotics}
\date{}

\begin{document}

\maketitle

\tableofcontents
\newpage

\section{引言}

在机器人学中，齐次变换矩阵$T(t) \in SE(3)$的时间导数与其twist $V$之间存在重要关系：

\begin{equation}
\dot{T} = [V] T
\label{eq:main_theorem}
\end{equation}

其中$[V]$是twist $V$的反对称矩阵表示（$4 \times 4$矩阵）。这个关系是理解刚体运动学和动力学的基础。本文将提供这个重要结果的详细证明。

\section{预备知识}

\subsection{齐次变换矩阵}

齐次变换矩阵$T \in SE(3)$的形式为：

\begin{equation}
T = \begin{bmatrix} R & p \\ 0 & 1 \end{bmatrix}
\end{equation}

其中：
\begin{itemize}
\item $R \in SO(3)$：旋转矩阵（$3 \times 3$）
\item $p \in \mathbb{R}^3$：平移向量
\item $SE(3)$：特殊欧氏群（Special Euclidean Group）
\end{itemize}

\subsection{Twist的定义}

\textbf{Twist} $V \in \mathbb{R}^6$定义为：

\begin{equation}
V = \begin{bmatrix} \omega \\ v \end{bmatrix}
\end{equation}

其中：
\begin{itemize}
\item $\omega \in \mathbb{R}^3$：角速度向量
\item $v \in \mathbb{R}^3$：线速度向量
\end{itemize}

\subsection{Twist的反对称矩阵表示}

对于twist $V = (\omega, v)$，其反对称矩阵表示为$4 \times 4$矩阵：

\begin{equation}
[V] = \begin{bmatrix}
[\omega] & v \\
0 & 0
\end{bmatrix} = \begin{bmatrix}
0 & -\omega_z & \omega_y & v_x \\
\omega_z & 0 & -\omega_x & v_y \\
-\omega_y & \omega_x & 0 & v_z \\
0 & 0 & 0 & 0
\end{bmatrix}
\label{eq:V_bracket}
\end{equation}

其中$[\omega]$是角速度$\omega$的$3 \times 3$反对称矩阵：

\begin{equation}
[\omega] = \begin{bmatrix}
0 & -\omega_z & \omega_y \\
\omega_z & 0 & -\omega_x \\
-\omega_y & \omega_x & 0
\end{bmatrix}
\end{equation}

\subsection{旋转矩阵的导数}

对于旋转矩阵$R(t) \in SO(3)$，其导数满足：

\begin{equation}
\dot{R} = [\omega_s] R = R [\omega_b]
\label{eq:R_dot}
\end{equation}

其中：
\begin{itemize}
\item $\omega_s$：空间角速度（在固定坐标系中表示）
\item $\omega_b$：体角速度（在旋转坐标系中表示）
\end{itemize}

\section{证明方法1：从点的运动出发}

\subsection{基本思路}

考虑固定在刚体上的点。如果刚体以twist $V$运动，则点的速度可以通过$\dot{T} = [V]T$来描述。

\subsection{详细证明}

\textbf{步骤1}：考虑固定在刚体上的点

设$q_0$是固定在刚体上的一个点（在刚体坐标系中的坐标）。在固定坐标系中，该点的位置为：

\begin{equation}
q(t) = T(t) q_0 = R(t) q_0 + p(t)
\end{equation}

\textbf{步骤2}：计算点的速度

对时间求导：

\begin{align}
\dot{q}(t) &= \dot{R}(t) q_0 + \dot{p}(t) \\
&= [\omega_s] R(t) q_0 + v_s
\end{align}

其中：
\begin{itemize}
\item $\omega_s$是空间角速度（在固定坐标系中表示）
\item $v_s$是空间线速度（在固定坐标系中表示）
\end{itemize}

\textbf{步骤3}：使用twist表示

定义空间twist $V_s = \begin{bmatrix} \omega_s \\ v_s \end{bmatrix}$，则：

\begin{align}
\dot{q}(t) &= [\omega_s] R(t) q_0 + v_s \\
&= \begin{bmatrix} [\omega_s] & v_s \\ 0 & 0 \end{bmatrix} \begin{bmatrix} R(t) q_0 \\ 1 \end{bmatrix} \\
&= [V_s] \begin{bmatrix} R(t) q_0 \\ 1 \end{bmatrix}
\end{align}

\textbf{步骤4}：扩展到齐次坐标

在齐次坐标中，$q_0 = \begin{bmatrix} q_{0,xyz} \\ 1 \end{bmatrix}$，$q = \begin{bmatrix} q_{xyz} \\ 1 \end{bmatrix} = T q_0$。

因此：
\begin{equation}
\dot{q} = [V_s] q = [V_s] T q_0
\end{equation}

\textbf{步骤5}：直接计算$\dot{q}$

另一方面，直接对$q = T q_0$求导：

\begin{equation}
\dot{q} = \dot{T} q_0
\end{equation}

\textbf{步骤6}：比较两种表达式}

由于$q_0$是任意的，我们得到：

\begin{equation}
\dot{T} q_0 = [V_s] T q_0
\end{equation}

对所有$q_0$成立，因此：

\begin{equation}
\dot{T} = [V_s] T
\label{eq:proof1_result}
\end{equation}

\textbf{证毕！}

\section{证明方法2：直接计算矩阵元素}

\subsection{基本思路}

直接计算$\dot{T}$和$[V]T$的每个元素，证明它们相等。

\subsection{详细证明}

\textbf{步骤1}：计算$\dot{T}$

对于$T = \begin{bmatrix} R & p \\ 0 & 1 \end{bmatrix}$：

\begin{equation}
\dot{T} = \begin{bmatrix} \dot{R} & \dot{p} \\ 0 & 0 \end{bmatrix}
\label{eq:T_dot_calc}
\end{equation}

\textbf{步骤2}：计算$[V]T$

设$V = \begin{bmatrix} \omega \\ v \end{bmatrix}$，则：

\begin{align}
[V]T &= \begin{bmatrix} [\omega] & v \\ 0 & 0 \end{bmatrix} \begin{bmatrix} R & p \\ 0 & 1 \end{bmatrix} \\
&= \begin{bmatrix} [\omega]R & [\omega]p + v \\ 0 & 0 \end{bmatrix}
\label{eq:V_bracket_T}
\end{align}

\textbf{步骤3}：使用旋转矩阵的导数关系}

从方程(\ref{eq:R_dot})，我们知道：

\begin{equation}
\dot{R} = [\omega_s] R
\end{equation}

其中$\omega_s$是空间角速度。

\textbf{步骤4}：确定线速度}

对于平移部分，我们需要确定$v$。

考虑刚体原点的运动。如果刚体原点在固定坐标系中的位置是$p(t)$，则其速度为$\dot{p}(t)$。

但是，由于刚体在旋转，原点的速度不仅包括平移速度，还包括由于旋转产生的速度。

\textbf{更精确的分析}：

对于固定在刚体上的点，其速度由两部分组成：
\begin{enumerate}
\item 由于刚体平移产生的速度：$v_s$
\item 由于刚体旋转产生的速度：$\omega_s \times (q - p_0)$，其中$p_0$是旋转中心
\end{enumerate}

对于刚体原点本身（$q = p$），如果旋转中心在原点，则：

\begin{equation}
\dot{p} = v_s + \omega_s \times p
\end{equation}

使用反对称矩阵：
\begin{equation}
\dot{p} = v_s + [\omega_s] p
\end{equation}

因此：
\begin{equation}
v_s = \dot{p} - [\omega_s] p
\end{equation}

\textbf{步骤5}：验证等式}

现在验证$\dot{T} = [V_s]T$：

\begin{align}
[V_s]T &= \begin{bmatrix} [\omega_s]R & [\omega_s]p + v_s \\ 0 & 0 \end{bmatrix} \\
&= \begin{bmatrix} [\omega_s]R & [\omega_s]p + (\dot{p} - [\omega_s]p) \\ 0 & 0 \end{bmatrix} \\
&= \begin{bmatrix} [\omega_s]R & \dot{p} \\ 0 & 0 \end{bmatrix} \\
&= \begin{bmatrix} \dot{R} & \dot{p} \\ 0 & 0 \end{bmatrix} \quad \text{（因为$\dot{R} = [\omega_s]R$）} \\
&= \dot{T}
\end{align}

\textbf{因此}：$\dot{T} = [V_s]T$，其中$V_s = \begin{bmatrix} \omega_s \\ \dot{p} - [\omega_s]p \end{bmatrix}$是空间twist。

\textbf{证毕！}

\section{证明方法3：使用指数映射}

\subsection{基本思路}

利用$SE(3)$的指数映射性质来证明。

\subsection{详细证明}

\textbf{步骤1}：指数映射

任意$T \in SE(3)$可以表示为：

\begin{equation}
T(t) = \exp([S]\theta(t))
\end{equation}

其中$S$是螺旋轴，$\theta$是运动参数。

\textbf{步骤2}：对时间求导}

\begin{align}
\dot{T}(t) &= \frac{d}{dt}[\exp([S]\theta(t))] \\
&= [S]\dot{\theta}(t) \exp([S]\theta(t)) \\
&= [S]\dot{\theta}(t) T(t)
\end{align}

\textbf{步骤3}：识别twist}

定义twist $V = S\dot{\theta}$，则：

\begin{equation}
\dot{T} = [V] T
\end{equation}

\textbf{证毕！}

\section{空间Twist与Body Twist}

\subsection{两种表示}

\textbf{空间Twist} $V_s$：
\begin{itemize}
\item 在固定坐标系中表示
\item 满足：$\dot{T} = [V_s] T$
\end{itemize}

\textbf{Body Twist} $V_b$：
\begin{itemize}
\item 在运动坐标系中表示
\item 满足：$T^{-1} \dot{T} = [V_b]$
\end{itemize}

\subsection{关系推导}

从$\dot{T} = [V_s] T$，左乘$T^{-1}$：

\begin{equation}
T^{-1} \dot{T} = T^{-1} [V_s] T
\end{equation}

\textbf{关键恒等式}：对于反对称矩阵和齐次变换矩阵，有：

\begin{equation}
T^{-1} [V_s] T = [\text{Ad}_{T^{-1}}(V_s)]
\end{equation}

其中$\text{Ad}_{T^{-1}}$是伴随变换。

更直接地，可以证明：

\begin{equation}
T^{-1} [V_s] T = [V_b]
\end{equation}

其中$V_b$是Body Twist。

\textbf{因此}：
\begin{equation}
T^{-1} \dot{T} = [V_b]
\end{equation}

\section{具体例子}

\subsection{例子1：纯平移}

考虑纯平移运动，$R = I$（单位矩阵），$p(t) = v_0 t$。

\textbf{计算$\dot{T}$}：
\begin{equation}
\dot{T} = \begin{bmatrix} 0 & v_0 \\ 0 & 0 \end{bmatrix}
\end{equation}

\textbf{计算$[V]T$}：

空间twist为$V_s = \begin{bmatrix} 0 \\ v_0 \end{bmatrix}$（无旋转，只有平移）。

\begin{align}
[V_s]T &= \begin{bmatrix} 0 & v_0 \\ 0 & 0 \end{bmatrix} \begin{bmatrix} I & p \\ 0 & 1 \end{bmatrix} \\
&= \begin{bmatrix} 0 & v_0 \\ 0 & 0 \end{bmatrix} = \dot{T} \quad \checkmark
\end{align}

\subsection{例子2：绕固定轴旋转}

考虑绕$z$轴旋转，旋转角度为$\theta(t)$，无平移。

\textbf{旋转矩阵}：
\begin{equation}
R(t) = \begin{bmatrix}
\cos\theta(t) & -\sin\theta(t) & 0 \\
\sin\theta(t) & \cos\theta(t) & 0 \\
0 & 0 & 1
\end{bmatrix}
\end{equation}

\textbf{计算$\dot{R}$}：
\begin{equation}
\dot{R}(t) = \dot{\theta}(t) \begin{bmatrix}
-\sin\theta(t) & -\cos\theta(t) & 0 \\
\cos\theta(t) & -\sin\theta(t) & 0 \\
0 & 0 & 0
\end{bmatrix}
\end{equation}

\textbf{空间角速度}：
\begin{equation}
\omega_s = \begin{bmatrix} 0 \\ 0 \\ \dot{\theta}(t) \end{bmatrix}
\end{equation}

\textbf{计算$[\omega_s]R$}：
\begin{align}
[\omega_s]R &= \begin{bmatrix}
0 & -\dot{\theta}(t) & 0 \\
\dot{\theta}(t) & 0 & 0 \\
0 & 0 & 0
\end{bmatrix} \begin{bmatrix}
\cos\theta(t) & -\sin\theta(t) & 0 \\
\sin\theta(t) & \cos\theta(t) & 0 \\
0 & 0 & 1
\end{bmatrix} \\
&= \dot{\theta}(t) \begin{bmatrix}
-\sin\theta(t) & -\cos\theta(t) & 0 \\
\cos\theta(t) & -\sin\theta(t) & 0 \\
0 & 0 & 0
\end{bmatrix} = \dot{R}(t) \quad \checkmark
\end{align}

\textbf{因此}：$\dot{T} = [V_s]T$，其中$V_s = \begin{bmatrix} \omega_s \\ 0 \end{bmatrix}$。

\section{应用}

\subsection{在速度递推中的应用}

在机器人动力学中，$\dot{T} = [V]T$用于速度递推：

\begin{equation}
V_i = \text{Ad}_{T_{i,i-1}}(V_{i-1}) + A_i \dot{\theta}_i
\end{equation}

这个递推关系的基础就是$\dot{T} = [V]T$。

\subsection{在Body雅可比计算中的应用}

从$T_{0i}^{-1}\dot{T}_{0i} = [V_i^b]$可以提取Body雅可比：

\begin{equation}
V_i^b = J_i^b(\theta) \dot{\theta}
\end{equation}

\section{总结}

\subsection{主要结论}

\begin{enumerate}
\item \textbf{基本关系}：$\dot{T} = [V_s]T$，其中$V_s$是空间twist

\item \textbf{Body形式}：$T^{-1}\dot{T} = [V_b]$，其中$V_b$是Body twist

\item \textbf{物理意义}：
\begin{itemize}
\item 空间twist $V_s$：从固定坐标系观察到的速度
\item Body twist $V_b$：从运动坐标系自身观察到的速度
\end{itemize}

\item \textbf{关系}：$V_s = \text{Ad}_T(V_b)$
\end{enumerate}

\subsection{证明方法总结}

\begin{itemize}
\item \textbf{方法1}：从点的运动出发，通过速度分析
\item \textbf{方法2}：直接计算矩阵元素，验证等式
\item \textbf{方法3}：使用指数映射的性质
\end{itemize}

所有方法都证明了相同的结果：$\dot{T} = [V]T$。

\section{参考文献}

\begin{itemize}
\item Modern Robotics: Mechanics, Planning, and Control
\item 第3章：刚体运动（Twist和Screw理论）
\item 第8章：开链动力学
\end{itemize}

\end{document}

